\documentclass[11pt]{article}
\usepackage[margin=1in]{geometry}
\usepackage{booktabs}
\usepackage{amsmath,amssymb}
\usepackage[hidelinks]{hyperref}
\setlength{\parskip}{4pt}
\title{\textbf{Pre-registered prediction: TESS Sector 114 transits of the 5.45-day candidate on TOI-913 (TIC 407126408)}}
\author{Josh Gold\\ \small Independent researcher, Acton, Massachusetts, USA}
\date{\small Registered 11 October 2026, before Sector 114 is observed (planned 2027 February 13 -- March 14).\\ Companion to doi:\href{https://doi.org/10.5281/zenodo.23290579}{10.5281/zenodo.23290579}}
\begin{document}
\sloppy
\maketitle

\section*{Purpose}
The candidate in the companion preprint has SNR $\approx 9$ over eight TESS sectors. Its credibility rests partly on choices made after seeing the data. This note fixes, in advance, the prediction for one new sector of TESS data, the measurement that will be made, and what each outcome will mean. It removes any room to adjust the analysis afterwards. We commit to publishing the result, whatever it is, as a new version of this record within 30 days of the Sector 114 light curves appearing on MAST.

\section*{Prediction}
Ephemeris from the joint fit in the companion preprint: $P = 5.450138 \pm 0.000016$\,d, $T_0 = 2458628.6479$\,BJD$_\mathrm{TDB}$ ($\sigma_{T_0}$ taken conservatively as 0.0073\,d), $T_{14} = 2.37$\,h. Uncertainties use $\sigma_n = \sqrt{\sigma_{T_0}^2 + (n\,\sigma_P)^2}$ and ignore the negative $T_0$--$P$ covariance, so they are conservative.

\begin{center}\footnotesize
\begin{tabular}{rcccc}
\toprule
Epoch & Mid-transit (BJD$_\mathrm{TDB}$) & $\sigma$ (min) & Mid-transit (UTC, geocentric) & Window (UTC) \\
\midrule
518 & 2461451.81938 & 15.9 & 2027-02-15 07:40 & 06:29--08:51 \\
519 & 2461457.26952 & 15.9 & 2027-02-20 18:28 & 17:17--19:39 \\
520 & 2461462.71966 & 15.9 & 2027-02-26 05:16 & 04:04--06:27 \\
521 & 2461468.16980 & 16.0 & 2027-03-03 16:03 & 14:52--17:14 \\
522 & 2461473.61994 & 16.0 & 2027-03-09 02:51 & 01:40--04:02 \\
523 & 2461479.07007 & 16.0 & 2027-03-14 13:39 & 12:28--14:50 (sector edge) \\
\bottomrule
\end{tabular}
\end{center}

Transits that fall in data gaps or outside the final sector boundaries will be reported and excluded. Predictions for 2026 October -- 2027 December, for ground-based observers, are in the accompanying CSV.

\section*{Frozen measurement}
\begin{itemize}\itemsep2pt
\item \textbf{Data:} SPOC 2-minute PDCSAP light curve for Sector 114 if available. Otherwise the TESS-SPOC full-frame-image light curve, which will be stated.
\item \textbf{Code:} tess-hunt commit \texttt{d7c59d1} (\url{https://github.com/killercatfish/tess-hunt}): the same preparation and detrending as the companion preprint, with TOI-913\,b's transits and the TOI-913.02 events masked.
\item \textbf{Statistic:} the ``core depth'' $D_{114}$, the mean depth over the central 60\% of each predicted transit relative to the local out-of-transit level. Its error $\sigma_{114}$ is computed from the scatter of out-of-transit windows in the same sector, as in the pipeline's sector-consistency test, with a 10\% systematic floor.
\item \textbf{Expected:} on the same measure, the eight existing sectors give $225 \pm 24$\,ppm, with typical single-sector errors of $\approx65$\,ppm. With the floor, we expect $\sigma_{114} \approx 69$\,ppm.
\end{itemize}

\section*{Pre-registered outcomes}
Using $D_\mathrm{exp} = 225$\,ppm and the measured $\sigma_{114}$:
\begin{center}\footnotesize
\begin{tabular}{lp{2.6in}cc}
\toprule
Outcome & Condition & $P$ if real & $P$ if no signal \\
\midrule
Insufficient data & fewer than 3 predicted transits with $\geq 50\%$ of the in-transit window covered by good data & --- & --- \\
Supports & $D_{114} > 2\sigma_{114}$ and $|D_{114} - 225| < 2\sigma_{114}$ & 0.90 & 0.02 \\
Inconclusive & otherwise, and not one of the rows below & 0.08 & 0.08 \\
Disfavours & $D_{114} < 225 - 2\sigma_{114}$ (i.e.\ $\lesssim 87$\,ppm) & 0.02 & 0.90 \\
Too deep & $D_{114} > 225 + 2\sigma_{114}$ & 0.02 & $\approx 0$ \\
\bottomrule
\end{tabular}
\end{center}
``Insufficient data'' is checked first; if it applies, no other outcome is reported and the test is deferred to the next sector that observes TOI-913. The probabilities assume $\sigma_{114} = 69$\,ppm and Gaussian errors. A ``too deep'' result would point to a blend or variability and will be reported as inconsistent with the planet interpretation as stated. Separately, the best-fit common time shift of the Sector 114 transits must lie within $\pm 48$\,min ($3\sigma$) of the prediction. A larger shift will be reported as an ephemeris failure.

After the pre-registered test, we will also report, as clearly labelled secondary results: the nine-sector refit, the updated SNR, and any changes to the vetting results. A single sector cannot validate the candidate. A ``supports'' outcome means only that the signal reappeared where and as predicted.

\section*{Integrity}
SHA-256 checksums of every file in this record are listed in \texttt{SHA256SUMS.txt}. The code is archived as GitHub release v0.1.0 of tess-hunt, which contains commit \texttt{d7c59d1}.

\section*{Acknowledgement}
The practice of pre-registering a held-out TESS sector follows P.~Rabtsevich (``A Pre-registered TESS Sector 110 Test of Two Transit-like Signals in TIC 4206066'', Zenodo, \url{https://doi.org/10.5281/zenodo.23175179}) and the \texttt{eumemic/exoplanets} project.
\end{document}
